\scanpage{7}

ent probably to clarify the very close relations one senses between the notions of $n$-category (and of $\infty$-category) on the one hand, that of simplicial set (or of topological space ...) on the other hand, just as the relation between $n$-Gr-categories (and $\unsure{6\!o|\infty}$-Gr-categories) and simplicial groups (or topological groups), finally between \deleted{\rlap{xxxxxxxxx}the notion} \deleted{\rlap{xx}of} Picard $n$-categories that are \uline{strict,} and \deleted{\rlap{xx}\unsure{the|the}} chain complexes truncated at dimension $n$. There would be reason to expect a \deleted{of}more or less complete fusion between these three visions: homological algebra, \deleted{\unsure{hat|hot|hom|ham}}homotopical\margin{\unsure{*|.|\phantom{*}}} algebra, categorical algebra, in a comm\uline{o}n vision (which might well take the name of \handwritten{``}noncommutative homological algebra\handwritten{''}, it being understood that what currently bears this name is only the beginning of what \deleted{remains to} \inserted{\handwritten{must}} come ...)
% The uncertain margin mark stands high beside the source line ending homóto-. Its phantom alternative records possible nontextual status; the physical mark is present.

As a last question, more down-to-earth, let us point out the one linked to the remark \deleted{that} \inserted{\handwritten{which follows.}} \deleted{w}\handwritten{W}hen \deleted{(for example)} $C$ is the category of finitely generated p\deleted{or}rojective modules over a \uline{commutative} ring $A$ (or more generally an additive category equipped with an \deleted{\unsure{associative, commutative, unital|associative, unital, commutative}} \inserted{\handwritten{associative, commutative, unital}} \uline{biadditive} operation $\otimes$), then the enveloping Picard category, in addition to its \deleted{str}obvious ``sum'' composition law (denoted here by $\oplus$ ) given \deleted{\rlap{xxxx}with} \inserted{by} its very definition, is equipped with an operation $\otimes$ coming fro\deleted{a}m the \inserted{\handwritten{ordinary}} tensor product operation of\handwritten{ the} modules, which also makes it an associative, commutative, unital $\otimes$-category (but obviously not a (Gr)-category), linked by obviou\deleted{e}\handwritten{s} distributivity isomorphisms \deleted{\handwritten{\unsure{to the|to}}} \deleted{\rlap{xxx}to the} \inserted{to the }operation\deleted{s} $\oplus$. It would be appropriate to study systematically the categories $A$ thus equipped with two operat\deleted{\rlap{x}a}ions $\deleted{\unsure{\otimes|\oplus}}\handwritten{\oplus}$ and $\otimes$, both associative, unital, commutative (additional data), with a distributivity datum
\[
X\otimes(Y\oplus Z)\handwritten{\simeq}X\otimes Y\oplus X\otimes Z
\]
and suitable compatibilities for the latter\deleted{\unsure{s}}; and to attempt to make a structure theory for such categories which are \deleted{\rlap{xxxxx}Picar} Picard for the $\oplus$ structure. This implies in particular that $\handwritten{\pi}_{0}(A)$ is a commutative unital ring, and $\handwritten{\pi}_{1}(A)$ a module over the latter. What other invariants must be introduced to characterize $A$ up to equivalence respecting all the structures ?

% archive sheet number: 6
